\section{Introduction}\label{sec:introduction}

For a superconducting accelerator cavity such as the nine-cell \SI{1.3}{GHz} niobium structure of the TESLA design \citep{aune2000superconducting,padamsee2009rf},
the designer needs the mode frequencies, the fields and the derived figures of merit:
the unloaded quality factor $\Qz$, the shunt impedance $\RoQ$, and the
peak-field ratios $\Epk/\Eacc$ and $\Bpk/\Eacc$ that limit the achievable gradient.
Each design iteration changes the shape and needs a new three-dimensional eigenvalue solve on a curved geometry.
Axisymmetric codes \citep{halbach1976superfish} are cheap but cover only part of the design space;
general three-dimensional design relies on finite-element eigensolvers such as the ACE3P suite \citep{ko2010advances} or commercial tools such as CST Studio Suite,
whose cost per geometry grows with the mesh resolution needed for accurate figures of merit.
A surrogate that returns frequencies, fields and figures of merit directly from the geometry would shorten the loop,
provided it can be trusted at the accuracy of the solver it replaces.

Neural operators \citep{lu2021learning,li2021fourier,kovachki2023neural}, which learn the map from a PDE's input to its solution, are natural candidates.
Applying them directly to the Maxwell eigenproblem $\curl\curl\,\E=k^2\E$ in $\Omega$ with $\mathbf{n}\times\E=0$ on $\partial\Omega$ is problematic for three reasons.
First, the network must return an ordered set of modes with mutually consistent eigenvalues and fields, and a plain regression loss guarantees none of this.
Second, the curl operator has an infinite-dimensional kernel formed by gradient fields;
any discretisation that does not respect the $H(\mathrm{curl})$ structure, and any predicted field that is not projected away from this kernel,
produces spurious modes at or near zero frequency \citep{hiptmair2002finite,boffi2010finite}.
Third, a network that outputs values on a fixed grid or degrees of freedom is tied to one discretisation, whereas eigenvalue accuracy depends on the mesh and the polynomial order.

\paragraph{Approach.}
We therefore do not ask the network for the modes themselves.
The network maps the geometry to a \emph{learned basis} of $m$ vector fields ($m$ is the basis width; $K\le m$ denotes the number of predicted modes).
A fixed linear map takes the basis to the degrees of freedom of a conforming high-order N\'ed\'elec space \citep{nedelec1980mixed} of order $p$ on a curved (isoparametric) mesh,
the basis is projected onto the discretely divergence-free subspace (the kernel projection, a float64 conjugate-gradient solve),
and a Rayleigh--Ritz step in the projected span returns the $K$ lowest frequencies and the fields.
The network is trained mainly to place the true modes inside the span of the basis; the Ritz step, not a regression head, computes frequencies and fields, and as a Galerkin projection it obeys min--max.
We implement it with NGSolve and Netgen \citep{schoberl1997netgen,schoberl2014cpp11}.

\paragraph{Contributions.}
\begin{enumerate}
  \item \textbf{A neural operator with a built-in discrete eigenvalue bound.}
  Every predicted eigenvalue bounds from above the corresponding nonzero discrete eigenvalue of the space in which the model runs,
  and the kernel projection excludes spurious zero modes (\cref{sec:method}, \cref{prop:upper-bound}).
  In the main held-out study every prediction stayed above the discrete eigenvalue of its space (smallest margin $f_{\mathrm{model}}/f_{\mathrm{FE}}-1=\SI{0.47}{\percent}$ over arms, geometries and modes); % src: docs/34 §8.2
  the bound is relative to the discrete, not the exact, eigenvalue and does not certify mode identity.
  % TODO pass 3: rank check of assumption (A) and Ritz-bound stress test with the a posteriori inexact-CG bound
  \item \textbf{Independence of the discretisation space.}
  The same weights run unchanged on a $p=3$ space and on a mesh refined once ($8\times$ tetrahedra); all predicted eigenvalues lay above the FE eigenvalue of the evaluated space. % src: docs/34 §4.2
  Accuracy does not improve there (frequency error up by \numrange{0.2}{0.7} percentage points), which we attribute, untested, to the learned basis. % src: docs/34 §4.2, referee E6
  \item \textbf{Single-sample representation capacity.}
  With quadrature-point tokens on a curved $p=2$ space, fitting one 179-tetrahedron sample to its own $p=2$ label reaches a relative eigenvalue error of \numrange{1e-6}{9e-6}. % src: docs/28 §3
  This shows the representation can express the modes, not that it generalises.
  \item \textbf{A data pipeline with analytic validation.}
  An adaptive $p=3$ label generator on curved meshes (\cref{sec:data}); on an analytic smooth pillbox (no pipes or fillets) it reproduces the exact values to \num{6e-7} in frequency, \num{3e-5} in $\Qz$, \num{1e-5} in $\RoQ$, \num{2e-4} in $\Epk$ and \num{4e-4} in $\Bpk$ at \num{157000} degrees of freedom. % src: docs/31 §1
  Removing repeated work cut an adaptive label at \num{450000} DOF from \SI{664}{s} to \SI{433}{s}. % src: docs/31 §3.8
  \item \textbf{Training on low-fidelity labels.}
  A low-fidelity (LF) label is one $p=2$ solve on the model mesh. The high-fidelity (HF) label used in all training experiments is one $p=3$ solve on the same start mesh, projected to $p=2$; the adaptive generator path is not exercised by these experiments (\cref{sec:multifidelity}). % src: referee E2, EDITOR_NOTES_2 decision 3
  On this mesh an LF label is $2.8$--$4.3\times$ cheaper (the range reflects machine load), % src: docs/34 §8.1, EDITOR_NOTES_2 decision 4
  differs from the HF label by \SIrange{0.004}{0.047}{\percent} in frequency and \SIrange{0.08}{0.32}{\percent} (mass norm) in the field subspace, % src: docs/34 §8.1
  and the model's errors measured against LF and against HF labels differ by only $0.02$ percentage points. % src: docs/34 §8.2
  In an equal-cost comparison on six held-out HF geometries, LF-trained and HF-trained models reached comparable accuracy (\cref{tab:equal-cost}).
  Limits: single runs; fidelity and the number of training geometries change together (no LF run on the HF geometries); and the label difference is one to two orders of magnitude smaller than the model error, so the study does not yet probe the regime where label type matters. % src: referee E1, EDITOR_NOTES_2 decision 2
  % TODO pass 3: repeated seeds / LF-8 control if produced; parametric-regression and coarse-FE baselines
\end{enumerate}

\paragraph{Scope.}
This is a method-and-pipeline paper with small-scale evidence: CPU training, at most 40 training geometries from one family (a pillbox with beam pipes), a 64-dimensional two-block trunk. % src: BRIEF §1; docs/34 §4, §8.2
The absolute frequency errors of \SIrange{1}{2}{\percent} reflect this under-trained model, not a limit of the method;
the figure-of-merit loss is off; fine-tuning on HF labels and mixed LF+HF training were inconclusive (\cref{sec:experiments}).
Single-threaded, the model path takes \SI{8.3}{s} (operators \SI{5.0}{s}, forward pass \SI{3.3}{s}) against \SI{18}{s} to assemble and solve the FE problem in the same space; % src: docs/34 §8.2, EDITOR_NOTES_2 decision 7
at these sizes this is about a factor of two, not a speed advantage.
The model cannot beat the FE solution of its own space, and no external baseline (parametric regression, coarse FE) is included yet (\cref{sec:discussion}). % TODO pass 3: baselines

\paragraph{Outline.}
\Cref{sec:related} reviews related work; \cref{sec:formulation} states the eigenproblem; \cref{sec:method} gives the model and the bound;
\cref{sec:data} and \cref{sec:multifidelity} cover the label generator and the two fidelities; \cref{sec:experiments} reports the experiments; \cref{sec:discussion} discusses limitations and future work.
